\documentclass[11pt]{article}

\usepackage[a4paper,margin=0.85in]{geometry}
\usepackage{amsmath,amssymb,amsthm}
\usepackage{booktabs}
\usepackage{longtable}
\usepackage{enumitem}
\usepackage{xcolor}
\usepackage{hyperref}
\usepackage{listings}
\usepackage{array}
\usepackage{titlesec}

\hypersetup{
  colorlinks=true,
  linkcolor=blue,
  urlcolor=blue,
  citecolor=blue
}

\definecolor{codebg}{RGB}{248,248,248}
\definecolor{codeframe}{RGB}{220,220,220}
\definecolor{codetext}{RGB}{35,35,35}

\lstdefinestyle{pythonstyle}{
  language=Python,
  basicstyle=\ttfamily\small\color{codetext},
  backgroundcolor=\color{codebg},
  frame=single,
  rulecolor=\color{codeframe},
  breaklines=true,
  showstringspaces=false,
  columns=fullflexible,
  keepspaces=true
}

\setlist[itemize]{leftmargin=1.3em}
\setlist[enumerate]{leftmargin=1.5em}
\titleformat{\section}{\Large\bfseries}{\thesection}{0.6em}{}
\titleformat{\subsection}{\large\bfseries}{\thesubsection}{0.6em}{}
\titleformat{\subsubsection}{\normalsize\bfseries}{\thesubsubsection}{0.6em}{}

\newcommand{\code}[1]{\texttt{#1}}
\newcommand{\given}{\,|\,}
\newcommand{\R}{\mathbb{R}}
\newcommand{\E}{\mathbb{E}}
\newcommand{\Var}{\operatorname{Var}}
\newcommand{\Cov}{\operatorname{Cov}}
\newcommand{\Normal}{\mathcal{N}}

\title{
  \textbf{MCMC Calibrator Interview Preparation Notes}\\
  \large Stamatics Society, IITK Project
}
\author{Prepared for technical interview defense}
\date{\today}

\begin{document}
\maketitle
\tableofcontents
\newpage

\section{How To Use This Document}

This document is written for an interview where the interviewer may ask questions
from quantitative finance, probability, stochastic processes, statistics,
algorithms, Python implementation, and software design. The goal is not to make
the project sound more advanced than it is. The goal is to know the exact
project deeply and to be able to defend every design choice from first
principles.

The current project is intentionally simple. It has one main script and three
small modules:

\begin{center}
\begin{tabular}{p{0.36\linewidth}p{0.54\linewidth}}
\toprule
\textbf{File} & \textbf{Purpose}\\
\midrule
\code{main.py} & Runs the whole workflow.\\
\code{data\_pipeline.py} & Downloads prices, computes hedge ratio, spread, and ADF test.\\
\code{ou\_model.py} & Computes Ornstein-Uhlenbeck log-likelihood.\\
\code{sampler.py} & Runs Metropolis-Hastings MCMC.\\
\bottomrule
\end{tabular}
\end{center}

If asked what the project does in one sentence:

\begin{quote}
This project takes two related financial assets, constructs a historical spread
between them, checks whether that spread is mean-reverting, and then uses a
Metropolis-Hastings MCMC sampler to estimate the hidden parameters of an
Ornstein-Uhlenbeck model governing that spread.
\end{quote}

\section{Project Elevator Pitch}

\subsection{The Problem}

In financial markets, individual stock prices often behave like noisy random
walks. Predicting one stock directly is hard because its movement depends on
news, macro factors, liquidity, investor behavior, and many unknown shocks.

Pairs trading changes the question. Instead of asking:

\begin{quote}
Where will stock A go tomorrow?
\end{quote}

we ask:

\begin{quote}
If stock A and stock B usually move together, has their relative price moved too
far away from its historical relationship?
\end{quote}

The object of interest becomes the \textbf{spread}, not the individual prices.

\subsection{The Core Idea}

Suppose there are two related assets, such as \code{GOOG} and \code{GOOGL}. If
their prices usually move together, then a suitable weighted difference between
their log-prices may fluctuate around a stable level. If that spread becomes too
high or too low, the model expects it to move back toward its average.

This property is called \textbf{mean reversion}.

\subsection{What We Estimate}

The spread is modeled as an Ornstein-Uhlenbeck process:

\[
dX_t = \theta(\mu - X_t)\,dt + \sigma\,dW_t.
\]

The unknown parameters are:

\begin{itemize}
  \item $\theta$: speed of mean reversion.
  \item $\mu$: long-term mean level of the spread.
  \item $\sigma$: volatility/noise level of the spread.
\end{itemize}

The project estimates these parameters using a Markov Chain Monte Carlo method,
specifically the Metropolis-Hastings algorithm.

\section{What The Project Is Not}

A smart interviewer may test whether you overclaim. Be precise:

\begin{itemize}
  \item It is not a live trading engine.
  \item It is not a guaranteed profitable strategy.
  \item It is not high-frequency trading.
  \item It does not place real orders.
  \item It does not include transaction costs, slippage, borrowing cost, or
  market impact.
  \item It does not prove the future will behave like the past.
  \item It is a research prototype for parameter calibration and understanding
  mean-reverting spreads.
\end{itemize}

Good interview answer:

\begin{quote}
The project demonstrates the statistical pipeline: construct a spread, test for
stationarity, specify an OU likelihood, and use MCMC to infer parameters. It is
not by itself a complete trading system because execution, costs, risk controls,
and regime monitoring are not implemented.
\end{quote}

\section{Day-Wise And Subtask-Wise Plan}

\subsection{Day 1: Understand The Trading Problem}

\textbf{Goal:} Understand why we use pairs rather than single-stock prediction.

\textbf{Subtasks:}

\begin{enumerate}
  \item Study statistical arbitrage and pairs trading.
  \item Understand that the spread is the main object, not the raw prices.
  \item Learn why mean reversion matters.
  \item Learn the difference between correlation and cointegration.
\end{enumerate}

\textbf{Interview questions:}

\begin{itemize}
  \item Why not predict stock price directly?
  \item Why use two assets?
  \item What does ``spread'' mean?
  \item What makes a pair tradable?
\end{itemize}

\textbf{Answer:} Direct price prediction is difficult because prices are often
non-stationary. A spread between related assets may be more stable and
mean-reverting. A pairs-trading strategy tries to exploit temporary deviations
from the historical relationship.

\subsection{Day 2: Data Collection With \code{yfinance}}

\textbf{Goal:} Download historical price data.

\textbf{Subtasks:}

\begin{enumerate}
  \item Use \code{yf.download()} for two tickers.
  \item Use adjusted close prices via \code{auto\_adjust=True}.
  \item Align dates by dropping missing values.
  \item Store prices in a Pandas \code{DataFrame}.
\end{enumerate}

\textbf{Why adjusted prices?} Adjusted prices account for corporate actions like
splits and dividends. Without adjustment, false jumps can enter the data.

\textbf{Interview trap:} If asked whether Yahoo Finance data is production-grade,
answer no. It is acceptable for research and student prototypes, but real
trading systems require reliable licensed data.

\subsection{Day 3: Log Prices And Hedge Ratio}

\textbf{Goal:} Estimate how much of asset B should be used to hedge asset A.

\textbf{Subtasks:}

\begin{enumerate}
  \item Convert raw prices to log prices.
  \item Regress $\log(A_t)$ on $\log(B_t)$.
  \item Extract intercept $\alpha$ and slope $\beta$.
  \item Interpret $\beta$ as the hedge ratio.
\end{enumerate}

\textbf{Regression model:}

\[
\log(A_t) = \alpha + \beta \log(B_t) + \epsilon_t.
\]

\textbf{Project spread definition:}

\[
X_t = \log(A_t) - \beta \log(B_t).
\]

The current code prints the intercept but does not subtract it from the spread.
That means the long-term mean $\mu$ estimated later absorbs the intercept-like
level of the spread.

\textbf{If asked whether one could also use residual spread:}

\[
\epsilon_t = \log(A_t) - \alpha - \beta \log(B_t).
\]

Answer: yes. That would center the spread more directly around zero. The current
project keeps the original notebook formula, where the OU mean $\mu$ handles the
level.

\subsection{Day 4: Stationarity And ADF Test}

\textbf{Goal:} Check whether the spread is suitable for mean-reversion modeling.

\textbf{Subtasks:}

\begin{enumerate}
  \item Run Augmented Dickey-Fuller test on the spread.
  \item Read the test statistic and p-value.
  \item Use p-value below 0.05 as a practical gate.
  \item Reject pairs whose spreads look non-stationary.
\end{enumerate}

\textbf{Interview answer:} A mean-reversion model only makes sense if the spread
does not wander away permanently. The ADF test checks evidence against the null
hypothesis of a unit root. If the p-value is low, we reject the unit-root null
and treat the spread as stationary for this prototype.

\subsection{Day 5: Ornstein-Uhlenbeck Model}

\textbf{Goal:} Choose a stochastic model for the spread.

\textbf{Subtasks:}

\begin{enumerate}
  \item Understand OU process.
  \item Derive conditional mean.
  \item Derive conditional variance.
  \item Write log-likelihood.
\end{enumerate}

\textbf{OU process:}

\[
dX_t = \theta(\mu - X_t)\,dt + \sigma dW_t.
\]

It has two forces:

\begin{itemize}
  \item deterministic pull back to $\mu$,
  \item random Brownian shock through $\sigma dW_t$.
\end{itemize}

\subsection{Day 6: Log-Likelihood}

\textbf{Goal:} Build a score function that says how plausible parameters are.

\textbf{Subtasks:}

\begin{enumerate}
  \item For each consecutive pair $(X_{t-1}, X_t)$, compute predicted mean.
  \item Compute predicted variance.
  \item Compute residual.
  \item Sum Gaussian log-probabilities.
\end{enumerate}

This is implemented in \code{ou\_model.py}.

\subsection{Day 7: MCMC And Metropolis-Hastings}

\textbf{Goal:} Estimate parameters by sampling possible values.

\textbf{Subtasks:}

\begin{enumerate}
  \item Start from initial guesses.
  \item Propose new parameters by random walk.
  \item Score old and new parameters.
  \item Accept or reject based on likelihood ratio.
  \item Store trace arrays.
\end{enumerate}

\subsection{Day 8: Burn-In And Parameter Extraction}

\textbf{Goal:} Convert raw MCMC traces into final values.

\textbf{Subtasks:}

\begin{enumerate}
  \item Discard first 25 percent of samples.
  \item Take medians of remaining samples.
  \item Report $\theta$, $\mu$, $\sigma$.
  \item Compute half-life.
\end{enumerate}

\subsection{Day 9: Interview Defense}

\textbf{Goal:} Know limitations and alternatives.

\textbf{Subtasks:}

\begin{enumerate}
  \item Explain why AMD/INTC failed.
  \item Explain why GOOG/GOOGL is a clean control pair.
  \item Explain why this is not guaranteed profit.
  \item Explain what you would add later, without claiming it is already built.
\end{enumerate}

\section{Code Flow From First Principles}

\subsection{Runtime Flow}

When we run:

\begin{lstlisting}[style=pythonstyle]
python main.py
\end{lstlisting}

the following happens:

\begin{enumerate}
  \item \code{main.py} imports NumPy, Matplotlib, \code{calculate\_hedge\_ratio},
  and \code{run\_mcmc\_calibrator}.
  \item It sets the tickers:
  \[
  A = \texttt{GOOG}, \quad B = \texttt{GOOGL}.
  \]
  \item It calls \code{calculate\_hedge\_ratio}.
  \item That function downloads historical prices.
  \item It computes log-prices.
  \item It runs OLS regression.
  \item It constructs the spread.
  \item It runs the ADF stationarity test.
  \item Control returns to \code{main.py}.
  \item \code{main.py} plots the spread.
  \item It converts the spread column to a NumPy array.
  \item It calls \code{run\_mcmc\_calibrator}.
  \item The sampler repeatedly calls \code{compute\_log\_likelihood}.
  \item It returns arrays for $\theta$, $\mu$, and $\sigma$.
  \item \code{main.py} discards burn-in and computes medians.
  \item It prints calibrated parameters and half-life.
\end{enumerate}

\subsection{Why The Code Is Split Into Modules}

This is not object-oriented in the heavy sense. It is simple functional
compartmentalization:

\begin{itemize}
  \item \code{data\_pipeline.py}: data and spread construction.
  \item \code{ou\_model.py}: mathematical likelihood.
  \item \code{sampler.py}: MCMC algorithm.
  \item \code{main.py}: orchestration.
\end{itemize}

\textbf{If asked why not use classes:} The current state does not require
persistent complex objects. Functions are enough because each phase takes input,
produces output, and does not need internal mutable object state. If the project
grew to support many pairs, configurations, result storage, and diagnostics,
classes like \code{PairData}, \code{OUSampler}, and \code{CalibrationResult}
could become useful.

\section{Financial Concepts}

\subsection{Asset}

An asset is a tradable financial instrument. In this project, assets are stocks
such as \code{GOOG} and \code{GOOGL}.

\subsection{Price Series}

A price series is a sequence:

\[
P_1, P_2, \dots, P_T
\]

where $P_t$ is the price at time $t$.

\subsection{Return}

A simple return is:

\[
R_t = \frac{P_t - P_{t-1}}{P_{t-1}}.
\]

A log return is:

\[
r_t = \log(P_t) - \log(P_{t-1}) = \log\left(\frac{P_t}{P_{t-1}}\right).
\]

The project does not directly model returns. It models a log-price spread.

\subsection{Why Log Prices?}

Log prices are useful because percentage changes become differences:

\[
\log(P_t) - \log(P_{t-1}) = \log(P_t/P_{t-1}).
\]

If one asset trades near 100 and another near 1000, raw price differences are
not scale-comparable. Log transformation makes relationships more proportional.

\subsection{Pairs Trading}

Pairs trading means taking positions in two assets simultaneously so that the
portfolio is less exposed to broad market movement and more exposed to relative
mispricing.

Example:

\begin{itemize}
  \item If spread is too high, asset A may be expensive relative to B.
  \item Strategy may short A and buy B.
  \item If spread falls back, the combined position profits.
\end{itemize}

This project does not implement the trade execution. It calibrates the
statistical model that could later inform such a strategy.

\subsection{Hedge Ratio}

The hedge ratio $\beta$ says how much of asset B is needed to offset asset A in
the spread:

\[
X_t = \log(A_t) - \beta \log(B_t).
\]

If $\beta \approx 1$, one unit of B roughly offsets one unit of A on log scale.

If $\beta > 1$, movements in B need a larger weight.

If $\beta < 0$, the assets move inversely in the regression relationship.

\subsection{Correlation Versus Cointegration}

Correlation measures whether two series move together at the same time. It is a
short-run co-movement measure.

Cointegration means a linear combination of two non-stationary series is
stationary. This is stronger and more relevant to pairs trading.

\textbf{Interview answer:} Two assets can be highly correlated but not
cointegrated. They may move in the same direction but still drift apart over
time. For mean-reversion trading, stationarity of the spread matters more than
simple correlation.

\section{Statistical Foundations}

\subsection{Random Variable}

A random variable is a quantity whose value depends on chance. In the project,
future spread values are treated as random variables.

\subsection{Probability Distribution}

A probability distribution describes how likely different values are. For a
continuous variable, we often use a density function.

\subsection{Expected Value}

The expected value is a probability-weighted average:

\[
\E[X] = \int x f(x)\,dx.
\]

In discrete form:

\[
\E[X] = \sum_i x_i p_i.
\]

\subsection{Variance}

Variance measures average squared deviation from the mean:

\[
\Var(X) = \E[(X - \E[X])^2].
\]

In this project, $\sigma$ controls the volatility/noise of the spread.

\subsection{Normal Distribution}

A normal distribution is written:

\[
X \sim \Normal(m, v),
\]

where $m$ is mean and $v$ is variance. Its density is:

\[
f(x) = \frac{1}{\sqrt{2\pi v}}
\exp\left(-\frac{(x-m)^2}{2v}\right).
\]

The OU transition uses a normal conditional distribution.

\section{OLS Regression From First Principles}

\subsection{The Model}

We model:

\[
y_t = \alpha + \beta x_t + \epsilon_t,
\]

where:

\begin{itemize}
  \item $y_t = \log(A_t)$,
  \item $x_t = \log(B_t)$,
  \item $\alpha$ is intercept,
  \item $\beta$ is slope or hedge ratio,
  \item $\epsilon_t$ is residual error.
\end{itemize}

\subsection{Objective}

OLS chooses $\alpha$ and $\beta$ to minimize squared residuals:

\[
\min_{\alpha,\beta}\sum_{t=1}^T
\left(y_t - \alpha - \beta x_t\right)^2.
\]

\subsection{Why Squared Errors?}

Squaring has useful properties:

\begin{itemize}
  \item positive and negative errors do not cancel,
  \item large errors are penalized more,
  \item the objective is differentiable,
  \item under Gaussian noise, OLS is maximum likelihood.
\end{itemize}

\subsection{Closed-Form Slope}

For simple regression:

\[
\hat{\beta}
= \frac{\sum_t (x_t - \bar{x})(y_t - \bar{y})}
{\sum_t (x_t - \bar{x})^2}
= \frac{\Cov(x,y)}{\Var(x)}.
\]

\[
\hat{\alpha} = \bar{y} - \hat{\beta}\bar{x}.
\]

\subsection{R-Squared}

\[
R^2 = 1 - \frac{\sum_t (y_t-\hat{y}_t)^2}
{\sum_t (y_t-\bar{y})^2}.
\]

It measures how much of the variation in $y$ is explained by the regression.

\textbf{Important:} High $R^2$ does not guarantee stationarity of the spread.
That is why we still need the ADF test.

\section{Stationarity And ADF}

\subsection{Stationarity}

A time series is stationary if its statistical behavior does not change over
time. Weak stationarity usually means:

\begin{itemize}
  \item constant mean,
  \item constant variance,
  \item autocovariance depends only on lag, not calendar time.
\end{itemize}

For pairs trading, stationarity means the spread is not drifting away
permanently.

\subsection{Unit Root}

A unit-root process has persistent shocks. A simple random walk is:

\[
X_t = X_{t-1} + \epsilon_t.
\]

If it goes up today, that shock stays embedded forever. Such a process does not
mean-revert.

\subsection{ADF Test}

The Augmented Dickey-Fuller test checks the null hypothesis:

\[
H_0: \text{series has a unit root}
\]

against:

\[
H_1: \text{series is stationary}
\]

If p-value $< 0.05$, we reject the unit-root null at the 5 percent level.

\subsection{ADF In The Code}

\begin{lstlisting}[style=pythonstyle]
adf_result = adfuller(data["Spread"])
\end{lstlisting}

The code prints:

\begin{itemize}
  \item ADF statistic,
  \item p-value,
  \item critical values,
  \item success or warning message.
\end{itemize}

\subsection{Interview Caution}

Do not say:

\begin{quote}
ADF proves this spread will always mean-revert.
\end{quote}

Say:

\begin{quote}
ADF gives statistical evidence against a unit root in the historical sample.
It supports using a mean-reversion model, but future regimes may change.
\end{quote}

\section{Stochastic Process Foundations}

\subsection{What Is A Stochastic Process?}

A stochastic process is a collection of random variables indexed by time:

\[
\{X_t : t \geq 0\}.
\]

In this project, $X_t$ is the spread at time $t$.

\subsection{Brownian Motion}

Brownian motion $W_t$ is a continuous-time random process with:

\begin{itemize}
  \item $W_0 = 0$,
  \item independent increments,
  \item $W_t - W_s \sim \Normal(0, t-s)$ for $t > s$,
  \item continuous paths.
\end{itemize}

The term $dW_t$ in the OU equation represents random market shock.

\subsection{Stochastic Differential Equation}

An SDE has deterministic and random parts:

\[
dX_t = a(X_t,t)\,dt + b(X_t,t)\,dW_t.
\]

For OU:

\[
a(X_t,t) = \theta(\mu-X_t),
\qquad
b(X_t,t)=\sigma.
\]

\section{Ornstein-Uhlenbeck Process}

\subsection{Definition}

\[
dX_t = \theta(\mu - X_t)\,dt + \sigma\,dW_t.
\]

\subsection{Parameter Meaning}

\begin{longtable}{p{0.18\linewidth}p{0.72\linewidth}}
\toprule
\textbf{Parameter} & \textbf{Meaning}\\
\midrule
$\mu$ & Long-run equilibrium level. The spread is pulled toward this.\\
$\theta$ & Pull strength or speed of mean reversion. Larger $\theta$ means faster reversion.\\
$\sigma$ & Noise intensity. Larger $\sigma$ means more random movement around the mean.\\
\bottomrule
\end{longtable}

\subsection{Why OU For This Project?}

OU is a natural model for mean-reverting continuous variables. The spread is
assumed to behave like a noisy rubber band: when far from $\mu$, the drift pulls
it back.

\subsection{Drift Direction}

The deterministic part is:

\[
\theta(\mu - X_t).
\]

If $X_t > \mu$, then $\mu-X_t < 0$, so drift is negative and pulls downward.

If $X_t < \mu$, then $\mu-X_t > 0$, so drift is positive and pulls upward.

If $X_t = \mu$, deterministic drift is zero.

\subsection{Exact Discrete Transition}

The code uses the exact OU transition:

\[
X_t \given X_{t-1}
\sim
\Normal(m_t, v_t),
\]

where:

\[
m_t = X_{t-1}e^{-\theta \Delta t}
+ \mu(1-e^{-\theta \Delta t}),
\]

\[
v_t =
\frac{\sigma^2}{2\theta}
\left(1-e^{-2\theta \Delta t}\right).
\]

In code:

\begin{lstlisting}[style=pythonstyle]
exp_term = np.exp(-theta * dt)
conditional_mean = x_old * exp_term + mu * (1.0 - exp_term)
conditional_var = (sigma**2 / (2.0 * theta)) * (
    1.0 - np.exp(-2.0 * theta * dt)
)
\end{lstlisting}

\subsection{Derivation Sketch}

Rewrite:

\[
dX_t = \theta\mu\,dt - \theta X_t\,dt + \sigma dW_t.
\]

Let:

\[
Y_t = X_t - \mu.
\]

Then:

\[
dY_t = -\theta Y_t\,dt + \sigma dW_t.
\]

Solving this linear SDE:

\[
Y_t = Y_s e^{-\theta(t-s)}
+ \sigma \int_s^t e^{-\theta(t-u)}\,dW_u.
\]

Thus:

\[
X_t = \mu + (X_s-\mu)e^{-\theta(t-s)}
+ \sigma \int_s^t e^{-\theta(t-u)}\,dW_u.
\]

The stochastic integral is normally distributed with mean zero and variance:

\[
\sigma^2\int_s^t e^{-2\theta(t-u)}\,du
=
\frac{\sigma^2}{2\theta}
\left(1-e^{-2\theta(t-s)}\right).
\]

This gives the transition used in the code.

\section{Likelihood And Log-Likelihood}

\subsection{What Is Likelihood?}

Probability asks:

\begin{quote}
Given parameters, how likely is the data?
\end{quote}

Likelihood asks:

\begin{quote}
Given observed data, how plausible are these parameters?
\end{quote}

For fixed data, likelihood is a function of parameters:

\[
L(\theta,\mu,\sigma \given X_1,\dots,X_T).
\]

\subsection{Conditional Likelihood}

Because the OU process is Markov:

\[
p(X_1,\dots,X_T)
= p(X_1)\prod_{t=2}^T p(X_t \given X_{t-1}).
\]

The code effectively scores:

\[
\prod_{t=2}^T p(X_t \given X_{t-1};\theta,\mu,\sigma).
\]

\subsection{Gaussian Log-Likelihood}

For:

\[
X_t \given X_{t-1} \sim \Normal(m_t,v_t),
\]

the density is:

\[
p(X_t\given X_{t-1})
=
\frac{1}{\sqrt{2\pi v_t}}
\exp\left(
-\frac{(X_t-m_t)^2}{2v_t}
\right).
\]

Taking log:

\[
\log p(X_t\given X_{t-1})
=
-\frac{1}{2}\log(2\pi v_t)
-\frac{(X_t-m_t)^2}{2v_t}.
\]

Summing over all transitions:

\[
\ell(\theta,\mu,\sigma)
=
-\frac{n}{2}\log(2\pi v)
-\frac{1}{2v}\sum_{t=2}^T (X_t-m_t)^2.
\]

This is exactly the expression in the code.

\subsection{Why Use Logs?}

Raw likelihood multiplies many small probabilities. This can underflow to zero
in floating-point arithmetic.

Logs convert products into sums:

\[
\log\left(\prod_i p_i\right) = \sum_i \log(p_i).
\]

This is more numerically stable.

\subsection{Invalid Parameters}

The code rejects:

\[
\theta \leq 0,\qquad \sigma \leq 0.
\]

Why?

\begin{itemize}
  \item $\theta \leq 0$ would mean no positive pull back to mean.
  \item $\sigma \leq 0$ is invalid because volatility cannot be negative.
\end{itemize}

Code:

\begin{lstlisting}[style=pythonstyle]
if theta <= 0 or sigma <= 0:
    return -np.inf
\end{lstlisting}

Returning $-\infty$ makes such proposals impossible to accept under the
likelihood comparison.

\section{MCMC From First Principles}

\subsection{Why MCMC?}

We want the posterior-like distribution of parameters:

\[
p(\theta,\mu,\sigma \given \text{data}).
\]

Instead of only finding one best parameter value, MCMC creates a sequence of
parameter samples. After burn-in, those samples represent plausible parameter
values according to the data and model.

\subsection{Markov Chain}

A Markov chain is a sequence:

\[
S_0,S_1,S_2,\dots
\]

where the next state depends only on the current state:

\[
P(S_{t+1}\given S_t,S_{t-1},\dots,S_0)
=
P(S_{t+1}\given S_t).
\]

In this project:

\[
S_i = (\theta_i,\mu_i,\sigma_i).
\]

The next parameter proposal depends only on the current parameter values.

\subsection{Monte Carlo}

Monte Carlo means using randomness to approximate a quantity. Here, random
sampling approximates the parameter distribution.

\subsection{Markov Chain Monte Carlo}

MCMC combines both:

\begin{itemize}
  \item Markov chain: each new sample depends only on current sample.
  \item Monte Carlo: randomness explores the parameter space.
\end{itemize}

\subsection{Metropolis-Hastings Algorithm}

At current state $s$:

\[
s = (\theta,\mu,\sigma).
\]

Propose:

\[
s' = s + \eta,
\]

where $\eta$ is Gaussian random noise.

Compute:

\[
\log r = \ell(s') - \ell(s).
\]

Accept if:

\[
\log u < \log r,
\qquad u\sim Uniform(0,1).
\]

If accepted, move to $s'$. If rejected, stay at $s$.

\subsection{Why Accept Worse Moves?}

If $s'$ has lower likelihood, $\log r < 0$. It may still be accepted with
probability:

\[
\exp(\ell(s')-\ell(s)).
\]

This prevents the chain from getting stuck too early in one region.

\subsection{Proposal Widths}

Current code:

\begin{lstlisting}[style=pythonstyle]
proposal_widths = {"theta": 0.005, "mu": 0.00025, "sigma": 0.00025}
\end{lstlisting}

These are standard deviations of proposal jumps.

If too large:

\begin{itemize}
  \item proposals jump into unlikely areas,
  \item acceptance rate becomes too low,
  \item chain gets stuck.
\end{itemize}

If too small:

\begin{itemize}
  \item proposals are accepted often,
  \item but movement is slow,
  \item chain explores poorly.
\end{itemize}

\subsection{Acceptance Rate}

\[
\text{Acceptance rate}
=
\frac{\text{accepted proposals}}{\text{total proposals}}.
\]

For random-walk Metropolis in multiple dimensions, a commonly cited theoretical
target is around 23 percent under idealized assumptions. In practical small
projects, a broad functional range such as 15--40 percent can be acceptable if
the traces mix and results are stable.

\textbf{Interview answer:} Acceptance rate is a diagnostic, not the final proof
of correctness. Trace plots, repeated runs, and sensitivity checks are also
important.

\section{Burn-In}

\subsection{What Is Burn-In?}

The chain starts from arbitrary initial values:

\[
\theta_0=1.0,\quad
\mu_0=\text{sample mean},\quad
\sigma_0=\text{sample std}.
\]

Early samples may reflect the starting point more than the target distribution.
Burn-in discards these early samples.

Current code:

\begin{lstlisting}[style=pythonstyle]
burn_in = int(0.25 * len(trace_theta))
clean_theta = trace_theta[burn_in:]
\end{lstlisting}

This discards first 25 percent.

\subsection{Why Median?}

The code uses:

\[
\hat{\theta} = \text{median of post-burn-in } \theta \text{ samples}.
\]

Median is robust to outliers. Mean could also be used.

\section{Half-Life}

\subsection{Definition}

For OU mean reversion, expected deviation decays like:

\[
\E[X_t-\mu \given X_0] = (X_0-\mu)e^{-\theta t}.
\]

Half-life $t_{1/2}$ is the time when expected deviation is half:

\[
e^{-\theta t_{1/2}} = \frac{1}{2}.
\]

Taking logs:

\[
-\theta t_{1/2} = \log(1/2) = -\log 2.
\]

Thus:

\[
t_{1/2} = \frac{\log 2}{\theta}.
\]

\subsection{Interpretation}

If half-life is 8.37 days, it does not mean the spread must halve exactly after
8.37 days. It means that under the fitted OU model, the expected deviation from
mean decays by half in that time.

\subsection{Trading Interpretation}

Half-life can be used as a time-based risk check:

\begin{itemize}
  \item If spread has not reverted after several half-lives, model assumptions
  may have broken.
  \item Shorter half-life means faster capital recycling.
  \item Extremely short half-life may be untradable after costs.
  \item Extremely long half-life may not be useful for a practical pairs trade.
\end{itemize}

\section{Algorithmic Complexity}

\subsection{Data Pipeline}

Let $T$ be number of time points.

\begin{itemize}
  \item Downloading data: external network/API cost.
  \item Log transform: $O(T)$.
  \item OLS for one regressor: effectively $O(T)$.
  \item ADF test: depends on lag selection, but roughly linear to low-order
  polynomial in $T$ for this small use case.
\end{itemize}

\subsection{Likelihood}

Each likelihood call uses vectorized operations over all $T-1$ transitions.

Cost:

\[
O(T).
\]

\subsection{MCMC}

Let $N$ be iterations.

Each iteration calls likelihood once, so total cost:

\[
O(NT).
\]

Memory:

\[
O(N)
\]

for each trace array. There are three arrays: $\theta$, $\mu$, $\sigma$.

\subsection{Why NumPy Vectorization Helps}

The likelihood avoids a Python loop over all data points. Instead, NumPy applies
operations over arrays in optimized C-level loops.

This is why:

\begin{lstlisting}[style=pythonstyle]
x_old = data[:-1]
x_new = data[1:]
\end{lstlisting}

is better than manually iterating through each adjacent pair in Python.

\section{Python And Implementation Details}

\subsection{\code{if \_\_name\_\_ == "\_\_main\_\_"}}

This block runs only when the file is executed directly:

\begin{lstlisting}[style=pythonstyle]
if __name__ == "__main__":
    ...
\end{lstlisting}

If the file is imported by another file, this block does not run.

\subsection{Pandas DataFrame}

A DataFrame is a table-like structure with rows and columns. Here, dates are the
index and ticker symbols are columns.

\subsection{NumPy Array}

NumPy arrays store numerical data efficiently and support vectorized operations.
The spread is converted using:

\begin{lstlisting}[style=pythonstyle]
spread_data = processed_df["Spread"].to_numpy()
\end{lstlisting}

\subsection{Slicing}

\begin{lstlisting}[style=pythonstyle]
x_old = data[:-1]
x_new = data[1:]
\end{lstlisting}

If:

\[
data = [X_0,X_1,X_2,X_3],
\]

then:

\[
x\_old=[X_0,X_1,X_2],
\qquad
x\_new=[X_1,X_2,X_3].
\]

This pairs each day with the next day.

\subsection{Why Preallocate Trace Arrays?}

The code uses:

\begin{lstlisting}[style=pythonstyle]
trace_theta = np.zeros(iterations)
\end{lstlisting}

Preallocating arrays avoids repeated dynamic resizing.

\section{OOP And Design Questions}

\subsection{Is This Object-Oriented?}

Not strongly. It is mainly procedural/functional.

\textbf{Good answer:} I kept the implementation simple because the project has
one pipeline and one model. The modules separate responsibilities. If the
project expanded, object-oriented design could help represent reusable concepts.

\subsection{Possible Class Design If Needed}

If asked how to convert it to OOP:

\begin{lstlisting}[style=pythonstyle]
class PairData:
    def fetch_prices(self): ...
    def compute_spread(self): ...

class OUModel:
    def log_likelihood(self, theta, mu, sigma): ...

class MCMCSampler:
    def run(self): ...

class CalibrationResult:
    def half_life(self): ...
\end{lstlisting}

\subsection{Why Not Use Classes Now?}

Because unnecessary classes would make the project look more complex without
adding real value. Functions are enough for:

\begin{itemize}
  \item one data pipeline,
  \item one likelihood,
  \item one sampler,
  \item one script.
\end{itemize}

\section{Limitations And Honest Defenses}

\subsection{GOOG/GOOGL As A Control Pair}

GOOG and GOOGL are share classes of the same company. This makes them a clean
control pair for testing the mathematical engine.

But:

\begin{itemize}
  \item margins may be tiny,
  \item institutional competition is strong,
  \item transaction costs can dominate,
  \item it is not necessarily profitable for retail execution.
\end{itemize}

\textbf{Interview answer:} I used GOOG/GOOGL because it gives a clean
stationary spread to validate the calibration workflow. It is a mathematical
sandbox, not a claim that this pair is the best live trade.

\subsection{AMD/INTC Failure}

AMD and Intel are in the same broad semiconductor sector, but during recent
market regimes their business performance diverged. Same sector does not imply
stationary spread.

\textbf{Lesson:} The model must test stationarity; it cannot assume it.

\subsection{No Transaction Costs}

Transaction costs include:

\begin{itemize}
  \item brokerage fees,
  \item bid-ask spread,
  \item slippage,
  \item short borrowing costs,
  \item taxes,
  \item market impact.
\end{itemize}

This project does not include them.

\subsection{No Out-Of-Sample Backtest}

The project calibrates parameters on historical data. A full trading project
would split data into training and testing periods.

\subsection{No Priors}

The current sampler uses likelihood and implicit hard constraints
($\theta>0,\sigma>0$). It does not explicitly include Bayesian priors.

If asked whether it is fully Bayesian:

\begin{quote}
It is MCMC-style parameter sampling using a likelihood and hard parameter
constraints. A more complete Bayesian version would multiply likelihood by
explicit priors for $\theta$, $\mu$, and $\sigma$.
\end{quote}

\section{Common Interview Questions And Strong Answers}

\subsection{Project-Level Questions}

\subsubsection{Q: What is the goal of this project?}

The goal is to build a simple MCMC-based calibrator for a mean-reverting
pairs-trading spread. It downloads historical prices, estimates a hedge ratio,
constructs a spread, checks stationarity, models the spread as an OU process,
and uses Metropolis-Hastings to estimate $\theta$, $\mu$, and $\sigma$.

\subsubsection{Q: Why did you choose pairs trading?}

Because individual prices are often non-stationary and hard to predict. A spread
between related assets can be more stable. Pairs trading focuses on relative
mispricing rather than absolute direction.

\subsubsection{Q: Why MCMC?}

MCMC gives a sampling-based way to explore plausible parameter values. Instead
of only returning one point estimate, it builds a trace of parameter values that
are likely under the model.

\subsubsection{Q: Is this a production trading system?}

No. It is a research/calibration prototype. Production would need robust data,
out-of-sample testing, transaction costs, execution logic, monitoring, and risk
controls.

\subsection{Data Questions}

\subsubsection{Q: Why \code{yfinance}?}

It is convenient and free for prototyping. It is not institutional-grade data.

\subsubsection{Q: Why drop NaN rows?}

Both assets must be aligned on the same dates. If one asset has missing data on
a date, the spread cannot be computed reliably for that date.

\subsubsection{Q: Why use adjusted prices?}

Adjusted prices reduce artificial jumps caused by dividends and splits.

\subsection{Regression Questions}

\subsubsection{Q: What is the hedge ratio?}

It is the regression slope $\beta$ in:

\[
\log(A_t)=\alpha+\beta\log(B_t)+\epsilon_t.
\]

It tells how strongly asset A moves with asset B and how to weight B in the
spread.

\subsubsection{Q: What if beta is negative?}

Then the assets have an inverse relationship in the regression. The spread
formula still works mathematically, but trade interpretation changes.

\subsubsection{Q: Does high R-squared mean the pair is tradable?}

No. High R-squared means good in-sample linear fit, but the spread may still be
non-stationary. We need the ADF test.

\subsection{ADF Questions}

\subsubsection{Q: What is the null hypothesis of ADF?}

The null is that the series has a unit root, meaning it is non-stationary.

\subsubsection{Q: What does p-value below 0.05 mean?}

It means we reject the unit-root null at the 5 percent level and treat the
spread as stationary for the purposes of this model.

\subsubsection{Q: Is p-value 0.049 very different from 0.051?}

Not practically. The 0.05 cutoff is conventional. A good researcher would also
inspect plots, test robustness, and check out-of-sample behavior.

\subsection{OU Questions}

\subsubsection{Q: Why is OU mean-reverting?}

Because the drift term is $\theta(\mu-X_t)$. If $X_t$ is above $\mu$, drift is
negative. If below $\mu$, drift is positive.

\subsubsection{Q: What happens if theta increases?}

The pull back to mean becomes faster. Half-life decreases because:

\[
t_{1/2}=\frac{\log 2}{\theta}.
\]

\subsubsection{Q: What happens if sigma increases?}

The spread becomes noisier. Even with mean reversion, random shocks can push it
farther from the mean.

\subsubsection{Q: Why is sigma positive?}

Volatility is a scale parameter. Negative volatility has no meaning in the
normal distribution variance.

\subsection{Likelihood Questions}

\subsubsection{Q: What does log-likelihood measure?}

It measures how plausible a parameter set is given the observed spread path
under the OU transition model.

\subsubsection{Q: Why not use probability directly?}

Multiplying many small probabilities can underflow to zero. Logs are numerically
stable.

\subsubsection{Q: Why is the residual squared?}

Because the Gaussian density penalizes squared deviations from conditional mean.
Large misses receive larger penalty.

\subsection{MCMC Questions}

\subsubsection{Q: What is Markov about your chain?}

The next parameter proposal depends only on the current parameter state, not the
full past history.

\subsubsection{Q: What is Monte Carlo about it?}

It uses random sampling to explore parameter space.

\subsubsection{Q: Why accept worse proposals?}

To avoid getting stuck in local regions and to correctly sample from the target
distribution.

\subsubsection{Q: What is acceptance rate?}

It is the fraction of proposed parameter moves accepted by the sampler.

\subsubsection{Q: What is a good acceptance rate?}

There is no universal number, but for random-walk Metropolis in multiple
dimensions, around 20--40 percent is often healthy; around 23 percent is a
classic theoretical reference under ideal assumptions. The current project's
15.7 percent was functional after tuning, but diagnostics should not rely on
that number alone.

\subsection{Algorithm Questions}

\subsubsection{Q: What is time complexity?}

If $T$ is data length and $N$ is MCMC iterations, each likelihood is $O(T)$ and
the sampler is $O(NT)$.

\subsubsection{Q: Why preallocate arrays?}

It avoids repeated memory allocation during the loop and stores the trace
efficiently.

\subsubsection{Q: Why vectorize likelihood?}

NumPy performs array operations in optimized low-level code, which is faster
than Python loops.

\subsection{Software Design Questions}

\subsubsection{Q: Why split into modules?}

To separate responsibilities:

\begin{itemize}
  \item data pipeline,
  \item OU model,
  \item sampler,
  \item orchestration.
\end{itemize}

\subsubsection{Q: Why not add more abstractions?}

The current project is small. Adding classes and frameworks would create
unnecessary complexity. The structure is simple but still readable.

\subsection{Critical Thinking Questions}

\subsubsection{Q: What is the biggest weakness of the project?}

The biggest weakness is that it calibrates on one historical sample without
out-of-sample testing or transaction-cost modeling. It proves the calibration
workflow, not profitability.

\subsubsection{Q: What would you add next?}

I would add out-of-sample validation, rolling-window calibration, transaction
costs, and trace diagnostics. I would keep them separate from the current simple
calibrator.

\subsubsection{Q: Why did you not force AMD/INTC to work?}

Because a model should reject bad data. If the spread is non-stationary, forcing
MCMC on it gives misleading parameters.

\section{Line-By-Line Code Explanation}

\subsection{\code{data\_pipeline.py}}

\begin{lstlisting}[style=pythonstyle]
raw_data = yf.download(
    [ticker_a, ticker_b],
    start=start_date,
    end=end_date,
    auto_adjust=True,
    progress=False,
)
\end{lstlisting}

This downloads historical adjusted price data.

\begin{lstlisting}[style=pythonstyle]
data[ticker_a] = raw_data["Close"][ticker_a]
data[ticker_b] = raw_data["Close"][ticker_b]
data = data.dropna()
\end{lstlisting}

This extracts synchronized close prices and removes missing rows.

\begin{lstlisting}[style=pythonstyle]
log_a = np.log(data[ticker_a])
log_b = np.log(data[ticker_b])
\end{lstlisting}

This transforms prices into log prices.

\begin{lstlisting}[style=pythonstyle]
X = sm.add_constant(log_b)
model = sm.OLS(log_a, X).fit()
\end{lstlisting}

This fits:

\[
\log(A_t)=\alpha+\beta\log(B_t)+\epsilon_t.
\]

\begin{lstlisting}[style=pythonstyle]
data["Spread"] = log_a - (beta * log_b)
\end{lstlisting}

This creates the project spread.

\begin{lstlisting}[style=pythonstyle]
adf_result = adfuller(data["Spread"])
\end{lstlisting}

This tests whether the spread has a unit root.

\subsection{\code{ou\_model.py}}

\begin{lstlisting}[style=pythonstyle]
if theta <= 0 or sigma <= 0:
    return -np.inf
\end{lstlisting}

Invalid parameters are rejected immediately.

\begin{lstlisting}[style=pythonstyle]
x_old = data[:-1]
x_new = data[1:]
\end{lstlisting}

This constructs adjacent time pairs.

\begin{lstlisting}[style=pythonstyle]
conditional_mean = x_old * exp_term + mu * (1.0 - exp_term)
\end{lstlisting}

This computes OU conditional mean.

\begin{lstlisting}[style=pythonstyle]
conditional_var = (sigma**2 / (2.0 * theta)) * (
    1.0 - np.exp(-2.0 * theta * dt)
)
\end{lstlisting}

This computes OU conditional variance.

\begin{lstlisting}[style=pythonstyle]
residuals = x_new - conditional_mean
sum_squared_residuals = np.sum(residuals**2)
\end{lstlisting}

This measures prediction errors.

\begin{lstlisting}[style=pythonstyle]
log_likelihood = -0.5 * n * np.log(2.0 * np.pi * conditional_var) - (
    sum_squared_residuals / (2.0 * conditional_var)
)
\end{lstlisting}

This is the Gaussian log-likelihood.

\subsection{\code{sampler.py}}

\begin{lstlisting}[style=pythonstyle]
current_theta = 1.0
current_mu = data.mean()
current_sigma = data.std()
\end{lstlisting}

These are starting guesses.

\begin{lstlisting}[style=pythonstyle]
proposed_theta = current_theta + np.random.normal(0, proposal_widths["theta"])
\end{lstlisting}

This proposes a new $\theta$ by random walk.

\begin{lstlisting}[style=pythonstyle]
log_acceptance_ratio = proposed_log_lik - current_log_lik
\end{lstlisting}

Because log likelihoods are used, likelihood ratio becomes subtraction.

\begin{lstlisting}[style=pythonstyle]
if np.log(np.random.uniform(0, 1)) < log_acceptance_ratio:
    current_theta = proposed_theta
\end{lstlisting}

This is the Metropolis-Hastings accept/reject rule.

\section{Possible Whiteboard Derivations}

\subsection{Derive Half-Life}

Start from expected OU deviation:

\[
\E[X_t-\mu] = (X_0-\mu)e^{-\theta t}.
\]

Set:

\[
\E[X_t-\mu] = \frac{1}{2}(X_0-\mu).
\]

Cancel:

\[
e^{-\theta t}=\frac{1}{2}.
\]

Take log:

\[
-\theta t = -\log 2.
\]

Therefore:

\[
t=\frac{\log 2}{\theta}.
\]

\subsection{Derive Metropolis Acceptance Rule}

For symmetric proposal:

\[
q(s'\given s)=q(s\given s').
\]

Acceptance probability:

\[
\alpha(s,s')=\min\left(1,\frac{\pi(s')}{\pi(s)}\right).
\]

If $\pi(s)\propto L(s)$:

\[
\alpha=\min\left(1,\frac{L(s')}{L(s)}\right).
\]

Taking logs:

\[
\log\frac{L(s')}{L(s)} = \ell(s')-\ell(s).
\]

The code accepts if:

\[
\log u < \ell(s')-\ell(s).
\]

\section{What To Say If Pressed Hard}

\subsection{Is Your MCMC Actually Bayesian?}

Strict answer:

\begin{quote}
The implementation uses an MCMC mechanism over a likelihood with hard parameter
constraints. It does not include explicit prior distributions, so it is closer
to likelihood-based MCMC exploration than a fully specified Bayesian posterior.
To make it fully Bayesian, I would define priors such as log-normal or gamma for
$\theta$ and $\sigma$, and normal for $\mu$, then sample from log-posterior =
log-likelihood + log-prior.
\end{quote}

\subsection{Why Use MCMC If OU Has MLE?}

Good answer:

\begin{quote}
For the simple OU model, one can estimate parameters through maximum likelihood
or related closed-form/discrete methods. I used MCMC because the project goal is
to understand sampling-based calibration and uncertainty over parameters, not
only point estimation. MCMC also generalizes when the likelihood becomes more
complex.
\end{quote}

\subsection{What If The Spread Is Not Stationary?}

\begin{quote}
Then the OU assumption is not appropriate. The calibrator may still output
numbers, but they should not be trusted. The correct response is to reject the
pair or change the time window.
\end{quote}

\subsection{What If Acceptance Rate Is 0.07 Percent?}

\begin{quote}
The proposal widths are too large relative to the scale of the target
distribution. Most proposed parameters land in low-likelihood or invalid
regions, so the chain barely moves. I would shrink proposal widths and inspect
trace behavior.
\end{quote}

\subsection{What If Acceptance Rate Is 90 Percent?}

\begin{quote}
Proposal widths are probably too small. The chain accepts many moves but moves
very slowly, producing high autocorrelation.
\end{quote}

\section{Defensible Current Results}

From the notebook context, the clean control run produced approximately:

\begin{center}
\begin{tabular}{ll}
\toprule
\textbf{Quantity} & \textbf{Approximate Value}\\
\midrule
Pair & GOOG/GOOGL\\
Hedge ratio $\beta$ & 0.997\\
Intercept $\alpha$ & 0.022\\
$R^2$ & 0.9999\\
ADF p-value & 0.0292\\
Acceptance rate & 15.70 percent\\
$\theta$ & 0.0829\\
$\mu$ & 0.0224\\
$\sigma$ & 0.0013\\
Half-life & 8.37 days\\
\bottomrule
\end{tabular}
\end{center}

Defensible interpretation:

\begin{itemize}
  \item The pair is structurally very close because both are Alphabet share
  classes.
  \item The spread passed the ADF gate in the tested sample.
  \item The estimated $\theta$ implies an expected half-life around 8 trading
  days.
  \item This is useful for demonstrating the calibrator, not proof of a
  profitable real strategy.
\end{itemize}

\section{Final Interview Checklist}

Before the interview, be able to explain:

\begin{enumerate}
  \item Why pairs trading uses spread instead of raw prices.
  \item Why log prices are used.
  \item What OLS estimates.
  \item What hedge ratio means.
  \item Why high $R^2$ is not enough.
  \item What stationarity means.
  \item What ADF tests.
  \item What OU process means.
  \item What $\theta$, $\mu$, and $\sigma$ mean.
  \item How the OU transition distribution is computed.
  \item How the log-likelihood formula is built.
  \item Why logs are used.
  \item What MCMC means.
  \item What Markov means.
  \item What Monte Carlo means.
  \item How Metropolis-Hastings accepts/rejects.
  \item What proposal widths do.
  \item What acceptance rate means.
  \item Why burn-in is removed.
  \item Why median is used.
  \item How half-life is derived.
  \item Why half-life is not a guarantee.
  \item Why GOOG/GOOGL is a control pair.
  \item Why AMD/INTC failed.
  \item What the project does not include.
  \item How the code is compartmentalized.
  \item Why the project avoids unnecessary classes.
  \item What you would add next without claiming it is already built.
\end{enumerate}

\section{Short Version To Recite}

If you have only one minute:

\begin{quote}
This project builds a simple MCMC calibrator for a pairs-trading spread. First,
I download adjusted historical prices for two assets and convert them to log
prices. I regress log price of asset A on log price of asset B to estimate the
hedge ratio, then construct the spread as $\log(A)-\beta\log(B)$. I run an ADF
test to check if this spread is stationary, because mean-reversion modeling only
makes sense if the spread does not behave like a random walk. Then I model the
spread using an Ornstein-Uhlenbeck process with parameters $\theta$, $\mu$, and
$\sigma$. The likelihood uses the exact OU conditional normal transition. Since
the parameters are unknown, I use a Metropolis-Hastings MCMC sampler: propose
new parameters, compute log-likelihood, accept better moves and sometimes worse
moves according to the likelihood ratio, and store the trace. After burn-in, I
take medians as calibrated estimates. Finally, I compute half-life as
$\log(2)/\theta$, which gives the expected timescale of mean reversion. The
project is a calibration prototype, not a complete trading system.
\end{quote}

\section{Deep-Dive Question Bank}

\subsection{Probability And Statistics}

\begin{enumerate}
  \item What is a random variable?
  \item What is a probability distribution?
  \item What is a density?
  \item What is expectation?
  \item What is variance?
  \item Why is Gaussian distribution used here?
  \item What is likelihood?
  \item What is log-likelihood?
  \item What is maximum likelihood?
  \item What is a p-value?
  \item What is the null hypothesis in ADF?
  \item What is stationarity?
  \item What is unit root?
  \item What is autocorrelation?
  \item Why can high correlation still fail for pairs trading?
\end{enumerate}

\subsection{Quant Finance}

\begin{enumerate}
  \item What is pairs trading?
  \item What is statistical arbitrage?
  \item What is hedge ratio?
  \item What is spread?
  \item Why use log spread?
  \item What is cointegration?
  \item Why transaction costs matter?
  \item Why is GOOG/GOOGL clean but not necessarily profitable?
  \item Why can sector pairs fail?
  \item What is regime shift?
\end{enumerate}

\subsection{Stochastic Calculus}

\begin{enumerate}
  \item What is Brownian motion?
  \item What is an SDE?
  \item What is drift?
  \item What is diffusion?
  \item Why is OU mean-reverting?
  \item Derive the OU conditional mean.
  \item Derive the OU conditional variance.
  \item Derive half-life.
\end{enumerate}

\subsection{Algorithms}

\begin{enumerate}
  \item What is MCMC?
  \item What is Markov property?
  \item What is Monte Carlo?
  \item What is Metropolis-Hastings?
  \item Why accept worse moves?
  \item What is proposal distribution?
  \item What are proposal widths?
  \item What is burn-in?
  \item What is acceptance rate?
  \item What is time complexity?
  \item What is memory complexity?
\end{enumerate}

\subsection{Python And Design}

\begin{enumerate}
  \item Why use Pandas?
  \item Why use NumPy?
  \item What does vectorization mean?
  \item What does slicing do?
  \item Why preallocate arrays?
  \item Why split files into modules?
  \item Why not use classes?
  \item How would you convert this into OOP?
  \item How would you test this?
  \item How would you make it production-ready?
\end{enumerate}

\section{Closing Position}

The strongest final position in an interview is:

\begin{quote}
I understand this as a clean calibration prototype. The value of the project is
not that it guarantees profit, but that it connects the full chain from data to
statistical validation to stochastic modeling to MCMC parameter estimation to a
human-interpretable risk metric. I can explain the math and the code, and I also
know what the project deliberately does not yet include.
\end{quote}

\end{document}
